\section*{Introduction}
\addcontentsline{toc}{section}{Introduction} % adds to toc


In statistical data analysis and computer vision, we are often interested in focusing only on the intrinsic geometry of the data, not on its location in the ambient space it lives in. In data analysis, understanding the intrinsic geometry can be used to capture nonlinear patterns in the data, and, as a result, it is sometimes possible to represent the most prominent features of the dataset in a lower-dimensional space. In computer vision, understanding the intrinsic geometry can be used for solving correspondence problems.

Isolating the data's position from its local connectivity is a non-trivial task. One way to do that is to turn the dataset of interest into a weighted graph using kernels that satisfy certain conditions. After normalising the weighted graph, we can embed it using functions called diffusion maps. Diffusion maps capture local connectivity in the data by analysing how random walks move between vertices in the graph.

The variant of diffusion maps that we study in this thesis was presented by Nadler, Lafon, Coifman and Kevrekidis in 2005 (see \cite{nips:dif_maps}). They proved that, after turning the dataset into a weighted graph via the Gaussian kernel and embedding it via diffusion maps, the diffusion distance between any two vertices on the graph is proportional to the Euclidean distance in the embedded space. In the paper, the authors do not analyse in detail what conditions the kernel must satisfy for the embedding to work as expected. They also prove the main theorem, but rely on auxiliary results which are often not stated. Our aim is to prove the theorem in a self-contained manner and to identify sufficient conditions on the kernel for the theorem to be applicable.

The thesis is organised as follows. Section \ref{sec:one} explains the basic terms used in the thesis and contains key results used in later proofs. Section \ref{sec:two} defines diffusion distance and diffusion maps, and proves that the diffusion distance is proportional to the Euclidean distance in the embedded space. We begin by providing a complete proof of the spectral theorem for real symmetric matrices, since it is a key theorem we use later to define diffusion maps. After that, we build on \cite{nips:dif_maps}, but we prove the theorem in a self-contained way. We follow the finite Markov-matrix formulation of \cite{nips:dif_maps} and make explicit a set of assumptions on the kernel function under which the proof goes through.
\begin{comment}
    We start by stating the sufficient conditions the kernel must satisfy for the main theorem to hold; in the original paper, they only prove that the theorem holds for the Gaussian kernel.
\end{comment}
We also modify the proof since some of the steps in \cite{nips:dif_maps} were slightly ambiguous.

Section \ref{sec:three} applies the theory to both synthetic and established datasets. We illustrate the key properties of diffusion maps and show how they can be used for solving correspondence problems. We also illustrate how using diffusion maps differs from classical linear dimensionality reduction techniques. We use principal component analysis (see \cite[Section~1.1]{PCA}) as an example, since it is a well-known linear method. We finish by applying diffusion maps to the correspondence problem on figures from the Dynamic FAUST dataset (see \cite{FAUST}). Since the dataset contains non-rigid transformations, the embeddings do not align as cleanly as in the synthetic examples. We nonetheless believe the example is valuable, both as a demonstration using real data and as motivation for the discussion of what conditions must be satisfied for diffusion maps to work to their full potential. The code used to generate plots and calculate the embeddings in Section \ref{sec:three} can be found on GitHub (see \cite{git}). 

Large language models by Google (see \cite{gemini}), OpenAI (see \cite{gpt}), and Anthropic (see \cite{claude}) were used during the writing of this thesis to generate alternative ideas, better understand materials used, find relevant literature and references, find mistakes in proofs, help write and debug the Python code for creating plots and calculations, and improve the readability of existing text. None of the content in the thesis was written by artificial intelligence, and its output was thoroughly checked.